\chapter{Rank-One Rigidity and Unitary Extension}
\label{ch:rankone_rigidity_unitary_extension}

This appendix records two finite-dimensional linear-algebra arguments used in
fixed-point and channel representations: scalar rigidity on rank-one rays and
unitary extension from equality of Gram matrices.

\section{Rank-one ray rigidity}

The following theorem supplies the scalarity step in
Theorem~\ref{thm:positive_retraction_one_matrix_factor}.

\begin{theorem}[Rank-one ray scalar rigidity]
    \label{thm:rankone_ray_scalar_rigidity}
    \lean{WolfProps.exists_eq_smul_id_of_maps_rankOne_to_span}
    \leanok
    Let $T:\MN{D}\to\MN{D}$ be complex-linear. Suppose that for every
    $v\in\C^D$ there is a scalar $c_v\in\C$ such that
    $T(\ketbra{v}{v})=c_v\ketbra{v}{v}$.
    Then there is a scalar $c\in\C$ such that $T=c\,\id$.
\end{theorem}

\begin{proof}
    \leanok
    Write $P_v=\ketbra{v}{v}$. If $D=0$, then $\MN{D}$ is the zero vector
    space, so the conclusion holds with $c=0$. If $D=1$, the single matrix unit
    is $P_{e_0}$. The hypothesis gives $T(P_{e_0})=c_{e_0}P_{e_0}$, and
    linearity gives $T=c_{e_0}\,\id$ on all of $\MN{1}$.

    It remains to consider $D\geq2$. Set $c=c_{e_0}$. For $i\ne j$, the identities
    \begin{align}
        P_{e_i+e_j}+P_{e_i-e_j}
        &=2P_{e_i}+2P_{e_j},
        \notag\\
        c_{e_i+e_j}P_{e_i+e_j}+c_{e_i-e_j}P_{e_i-e_j}
        &=2c_{e_i}P_{e_i}+2c_{e_j}P_{e_j}
        \notag
    \end{align}
    give, from their $(i,i)$, $(j,j)$, and $(i,j)$ entries,
    $c_{e_i+e_j}=c_{e_i-e_j}=c_{e_i}=c_{e_j}$.
    The imaginary parallelogram identity and its image under $T$ are
    \begin{align}
        P_{e_i+\mathrm{i}e_j}+P_{e_i-\mathrm{i}e_j}
        &=2P_{e_i}+2P_{e_j},
        \notag\\
        c_{e_i+\mathrm{i}e_j}P_{e_i+\mathrm{i}e_j}
        +c_{e_i-\mathrm{i}e_j}P_{e_i-\mathrm{i}e_j}
        &=2c_{e_i}P_{e_i}+2c_{e_j}P_{e_j}.
        \notag
    \end{align}
    Comparing the same three entries gives
    $c_{e_i+\mathrm{i}e_j}=c_{e_i-\mathrm{i}e_j}=c_{e_i}=c_{e_j}$.
    Taking the pair $(0,i)$ shows that $c_{e_i}=c$ for every $i$.
    The hypothesis therefore gives $T(E_{ii})=cE_{ii}$ for the diagonal matrix
    units. For $i\ne j$, the polarization identity
    \begin{align}
        4E_{ij}
        &=P_{e_i+e_j}-P_{e_i-e_j}
        +\mathrm{i}P_{e_i+\mathrm{i}e_j}
        -\mathrm{i}P_{e_i-\mathrm{i}e_j}
        \notag
    \end{align}
    gives $T(E_{ij})=cE_{ij}$. Thus $T=c\,\id$ on every matrix unit.
\end{proof}

Thus rank-one domination in
Theorem~\ref{thm:positive_retraction_one_matrix_factor} produces one scalar
functional rather than a vector-dependent family of coefficients.

\section{Gram matrices and unitary extension}

The next two results supply the ambient unitary in the fixed-point decomposition
of Theorem~\ref{thm:wolf_6_14}.

\begin{lemma}[Rectangular Gram equality gives a unitary]
    \label{lem:exists_unitary_mul_eq_of_gram_eq}
    \lean{Matrix.exists_unitary_mul_eq_of_conjTranspose_mul_eq}
    \leanok
    \uses{}
    Let $A,B:\C^n\to\C^m$ be linear maps with
    $B^\dagger B=A^\dagger A$. Then there is a unitary
    $U$ on $\C^m$ such that $B=UA$.
\end{lemma}

\begin{proof}
    \leanok
    Let $\mathcal R_A=\operatorname{ran}A$ and
    $\mathcal R_B=\operatorname{ran}B$. Define
    $V:\mathcal R_A\to\mathcal R_B$ by $V(Ax)=Bx$. This is well defined:
    if $Ax=Ay$, then
    \begin{align}
        \lVert B(x-y)\rVert^2
        &=\langle x-y,B^\dagger B(x-y)\rangle
          =\langle x-y,A^\dagger A(x-y)\rangle
          =\lVert A(x-y)\rVert^2=0.
        \notag
    \end{align}
    The same calculation, with two vectors and polarization, gives
    $\langle V(Ax),V(Ay)\rangle=\langle Ax,Ay\rangle$. Thus $V$ is an
    isometry onto $\mathcal R_B$. In particular,
    $\dim\mathcal R_A=\dim\mathcal R_B$, so their orthogonal complements
    in $\C^m$ have equal dimension. Choose a unitary
    $V^\perp:\mathcal R_A^\perp\to\mathcal R_B^\perp$ and set
    $U=V\oplus V^\perp$ with respect to the two orthogonal decompositions of
    $\C^m$. Then $U$ is unitary and $UAx=Bx$ for every $x\in\C^n$; hence
    $B=UA$.
\end{proof}

\begin{theorem}[Unitary completion of an isometric inclusion]
    \label{thm:isometric_inclusion_zero_extension}
    \lean{Matrix.exists_unitary_zero_extension_eq}
    \leanok
    \uses{lem:exists_unitary_mul_eq_of_gram_eq}
    Let $W:\C^n\to\C^D$ satisfy $W^\dagger W=\Id_n$. There are an
    identification $\C^D\simeq\C^{D-n}\oplus\C^n$ and a unitary $U$ on
    $\C^D$ such that, for every $A\in\MN{n}$,
    \begin{align}
        WAW^\dagger
        &=U(0_{D-n}\oplus A)U^\dagger.
        \label{eq:representations_unitary_completion}
    \end{align}
\end{theorem}

\begin{proof}
    \leanok
    \uses{lem:exists_unitary_mul_eq_of_gram_eq}
    The isometry $W$ implies $n\leq D$. Let $J:\C^n\to
    \C^{D-n}\oplus\C^n$ be the canonical inclusion into the second summand.
    Since $W^\dagger W=J^\dagger J=\Id_n$,
    Lemma~\ref{lem:exists_unitary_mul_eq_of_gram_eq}
    gives a unitary $U$ with $W=UJ$. Substituting this identity gives
    (\ref{eq:representations_unitary_completion}) because
    $JAJ^\dagger=0_{D-n}\oplus A$.
\end{proof}

Equality of the two inclusion Gram matrices therefore supplies exactly the
unitary conjugation used to adjoin the complementary zero block in
Theorem~\ref{thm:wolf_6_14}.